\begin{recipe}{Nasi goreng}
\kicker[Hoofdgerecht,Vlees,Indonesisch,Doordeweeks]{Hoofdgerecht · vlees · Indonesisch · doordeweeks}
\index[register]{Nasi goreng}
\index[register]{Kip}
\index[register]{Rijst}

\meta{25 minuten}

\lettrine{N}{asi} goreng met kip, een gebakken ei en pindasaus. Dit recept gebruikt
een boemboe uit een pakje, wat het een stuk sneller maakt.

\blockrule{Ingrediënten}
\begin{ingredients}
  \ing{2}{kipfilets, in blokjes}\index[register]{Kip}
  \ing{5}{eieren (1 per persoon)}
  \ing{170 g}{langkorrelige rijst (b.v.\ Basmati)}\index[register]{Rijst}
  \ing{450 g}{nasi-bamigroenten}
  \ing{95 g}{nasi boemboe (b.v.\ Lidl)}
  \ing{50 ml}{water}
  \ing{60 g}{Calve pindasaus mild}
  \ing{120 ml}{melk voor pindasaus}
  \ingb{kokosolie}
  \ingb{gebakken uitjes}
  \ingb{kroepoek}
  \ingb{sambal, voor de liefhebber}
\end{ingredients}

\blockrule{Bereiding}
\tip{Rijst van de dag ervoor werkt het best. Koude rijst bakt droger en minder
plakkerig.}
\begin{steps}
  \item Kook water voor de rijst.
  \item Kook de rijst in 10 minuten, in het begin met deksel. Haal op tijd eraf voor overkoken. Roer ook soms even door zodat het niet aan de bodem plakt.
        Giet dan rijst af en laat zonder deksel afkoelen, zodat wat vocht eruit dampt.
  \item Terwijl rijst kookt, snij de kip in blokjes.
  \item Verhit een wokpan goed voor en doe een scheut kokosolie erin. Bak de
        kipblokjes op hoog vuur tot er aan alle kanten wat bruin op zit.
  \item Voeg de groenten toe en bak 4 minuten mee.
  \item Voeg de boemboe en 50 ml water toe en schep goed door.
  \item Voeg de rijst toe en roer op laag vuur door tot alles goed gemengd en warm is.
  \item Bak ondertussen een ei per persoon op middelhoog vuur in een koekenpan met ruim kokosolie en breng op smaak met peper en zout. Prik de dooiers door.
  \item Maak ondertussen in een klein sauspannetje de pindasaus met de melk op een
        laag vuur warm tot een gladde saus.
  \item Serveer met een gebakken ei erop en gebakken uitjes, kroepoek en sambal
        erbij. Wat zilveruitjes of komkommer ernaast is ook lekker.
\end{steps}

\heroimagefade{images/nasi}
\end{recipe}
